% ================================================================
%  Topological EWS Manuscript — medRxiv preprint
%  Compile: pdflatex euler_ews_medrxiv.tex (twice for references)
% ================================================================
\documentclass[12pt,a4paper]{article}

\usepackage[margin=2.5cm]{geometry}
\usepackage{amsmath,amssymb}
\usepackage{booktabs}
\usepackage{graphicx}
\usepackage{microtype}
\usepackage{setspace}
\usepackage{hyperref}
\usepackage[numbers,sort&compress]{natbib}
\usepackage{array}
\usepackage{caption}
\usepackage{xcolor}
\usepackage{rotating}
\usepackage{lineno}

\doublespacing
\linenumbers

\hypersetup{
  colorlinks=true,
  linkcolor=blue!60!black,
  citecolor=blue!60!black,
  urlcolor=blue!60!black
}

% ── title block ──────────────────────────────────────────────────
\title{\textbf{Heart rate oscillatory complexity is associated with septic
shock in two matched ICU cohorts}}

\author{
  Chengyi Fang$^{1}$ \and
  Haiping Yu$^{2,*}$ \and
  Qiancheng Luo$^{3,*}$
}

\date{
  $^1$Tongji University School of Medicine, Shanghai, China \\
  $^2$Department of Nursing, Pudong Gongli Hospital,
       Shanghai University of Medicine \& Health Sciences, Shanghai, China \\
  $^3$Department of Critical Care Medicine, Pudong Gongli Hospital,
       Shanghai University of Medicine \& Health Sciences, Shanghai, China \\[6pt]
  $^*$Co-corresponding authors \\[6pt]
  \textbf{Correspondence:}
  pingping670@sina.com (H.Y.); luoqiancheng19@163.com (Q.L.) \\[6pt]
  \textit{Preprint deposited on medRxiv. Not peer-reviewed.} \\[4pt]
  \textbf{Keywords:} Euler characteristic, topological data analysis,
  septic shock, early warning signal, heart rate variability,
  MIMIC-IV, eICU-CRD, sample entropy, oscillatory complexity
}

\begin{document}
\maketitle

% ================================================================
\begin{abstract}
Septic shock carries in-hospital mortality exceeding 40\%, yet
early-warning tools do not capture the loss of cardiovascular oscillatory
flexibility preceding vasopressor initiation.
We studied two topological heart rate (HR) complexity indicators---Euler
characteristic ($\chi_{\text{HR}}(0)$) and local-extrema count
($N_{\text{extrema}}$)---in a 1:2 matched MIMIC-IV cohort (725 shock,
1,447 controls) with external validation in eICU-CRD (587 pairs).
Both indicators declined in shock patients from 24\,h before vasopressor
initiation.
$N_{\text{extrema}}$ was associated with shock
(OR\,=\,0.88, 95\%\,CI 0.82--0.94; $p<0.001$), unchanged after adjusting
for concurrent mean HR and MAP (OR\,=\,0.87; $p<0.001$), and replicated
in eICU-CRD at $T_0\geq 72$\,h (OR\,=\,0.80, 0.65--0.98; $p=0.034$).
Lag-1 autocorrelation lost significance after mean adjustment
(OR\,=\,1.23; $p=0.271$), and sample entropy was non-significant without
adjustment (OR\,=\,0.86; $p=0.178$).
$N_{\text{extrema}}$ is mean-level-robust and replicable across ICU
settings; $\chi_{\text{HR}}(0)$ shows a supportive but partially
attenuated association after mean adjustment.
Both indicators are parameter-free and computable from routine hourly
EHR charting.
\end{abstract}

% ================================================================
\section{Introduction}
\label{sec:intro}

Septic shock complicates one in three ICU sepsis admissions and carries
in-hospital mortality exceeding 40\% despite modern management
\cite{singer2016}.
Existing severity scores (SOFA, APACHE) and threshold-based vital-sign
alerts quantify haemodynamic derangement but do not measure
cardiovascular \emph{regulatory capacity}---the oscillatory flexibility
of the heart rate (HR) that erodes in the hours before vasopressor
initiation \cite{scheffer2009}.
Capturing this erosion could enable earlier intervention in a condition
where every hour of delayed treatment worsens outcome.

Critical slowing down (CSD) theory predicts that a system near collapse
recovers more slowly from small perturbations, producing a signature of
rising lag-1 autocorrelation (AC1) and increasing variance
\cite{scheffer2009,dakos2012}.
Applied to hourly ICU vital signs, CSD-derived AC1 has shown associations
with clinical deterioration \cite{olderikkert2016,khalil2025}, yet its
independence from concurrent mean-level changes has not been rigorously
tested.
Entropy-based heart rate variability (HRV) metrics---sample entropy (SampEn)
and multiscale entropy (MSE)---quantify oscillatory irregularity and carry
prognostic information in critical illness \cite{moorman2011,lake2002,bravi2012},
but they require careful parameter selection (template length $m$, tolerance $r$)
and produce high rates of undefined values when applied to the sparse,
hourly EHR time series that are the standard of routine clinical monitoring.

Topological data analysis (TDA) characterises the \emph{shape} of a signal
through algebraic invariants that are robust to monotone rescaling and
require no free parameters \cite{edelsbrunner2010}.
For a locally detrended HR residual series, the Euler characteristic of
the sublevel-set filtration at threshold~0---denoted $\chi_{\text{HR}}(0)$---
counts the number of connected below-mean segments; it is equivalent to the
number of oscillatory excursions through the local mean and is computable in
$O(n)$ from as few as six observations per window.
The related local-extrema count ($N_{\text{extrema}}$) tallies directional
reversals without any threshold assumption, making it structurally
robust to mean-level shifts by construction.
To our knowledge, neither metric has been applied to ICU vital-sign
early-warning signals.

Here we report a matched-cohort study with three aims: (i)~to determine
whether $\chi_{\text{HR}}(0)$ and $N_{\text{extrema}}$ are associated with
late-onset septic shock in a large MIMIC-IV discovery cohort; (ii)~to test
whether these associations survive rigorous adjustment for concurrent mean
HR and MAP---the critical test of mean-level confounding independence;
and (iii)~to compare topological indicators against lag-1 autocorrelation
and sample entropy using parallel comparator analyses within the same
modeling framework, with external validation in eICU-CRD.

% ================================================================
\section{Methods}
\label{sec:methods}

\subsection{Data source and ethical considerations}
MIMIC-IV version 3.1 \cite{johnson2023} was accessed under PhysioNet Data
Use Agreement (GDOC-2010-0024).
All ICU stays within MIMIC-IV are de-identified; no additional ethical approval
was required.
Derived concepts (Sepsis-3, SOFA, vasopressor equivalents) were generated
using the MIT \texttt{mimic-code} DuckDB pipeline \cite{johnson2023}.

\subsection{Cohort construction}
\label{sec:cohort}
Full cohort construction details are reported in Fang et al.\ \cite{fang2026};
we briefly summarise the key steps.

From 94,458 MIMIC-IV ICU stays, 41,296 met Sepsis-3 criteria
\cite{singer2016}.
An operational septic shock definition was applied:
vasopressor initiation with lactate $>$2\,mmol/L ($\pm$24\,h),
vasopressor duration~$\geq$1\,h, and isotonic crystalloid~$\geq$1\,L in the
preceding 24\,h, yielding 10,310 shock stays.
The vasopressor start time was designated $T_0$.

Risk-set matching (1:2 ratio, admission SOFA\,$\pm$2, seed 42) produced 725
shock and 1,450 control (stay,\,$T_0$) pairs.
Controls were drawn from a dynamic risk set of Sepsis-3 stays with
ICU duration~$\geq H$\,h whose own shock onset (if any) fell after~$H$.
After time-series cleaning (see below), 2,172 (stay,\,$T_0$) observations
remained analysable (725 shock, 1,447 control) across 725 matched strata.

\subsection{Time-series preprocessing}
Hourly MAP (ABPm preferred over NBPm) and HR were extracted for
72\,h before each $T_0$.
Preprocessing followed the pipeline of Fang et al.\ \cite{fang2026}:
physiological-range filtering, 1-hour median resampling, causal 24-hour
trailing-mean detrending (the first 24\,h serving as burn-in only),
$\pm3\sigma$ outlier replacement, and exclusion of stays with
$<$24 actual observations in the final 48-hour analysis window.

\subsection{Topological indicators}
\label{sec:topo}

Two topological indicators of HR regulatory complexity were computed over
12-hour rolling windows (1-hour step) of the HR residual series,
using only non-interpolated, paired observations.
Windows with fewer than 6 actual observations were excluded.

\paragraph{Euler characteristic $\chi_{\text{HR}}(0)$.}
Let $r_1, r_2, \ldots, r_n$ denote the HR residual values (in temporal order)
at non-interpolated time points within a 12-hour window.
The sublevel set at threshold~0 is the discrete index set
$S_0 = \{i : r_i \leq 0\}$.
Its Euler characteristic equals the number of connected components of $S_0$:
\begin{equation}
  \chi_{\text{HR}}(0)
  = \mathbf{1}[r_1 \leq 0]
    + \sum_{i=2}^{n} \mathbf{1}[r_{i-1} > 0 \text{ and } r_i \leq 0].
  \label{eq:euler}
\end{equation}
This quantity counts the number of times HR crosses below its local trend
in each window.
A value of~0 indicates HR remains persistently above its trend (unrelenting
tachycardia without corrective return); a value of~1 indicates HR is
persistently below (bradycardic drift); higher values indicate frequent
oscillation through the trend level, reflecting preserved homeostatic
flexibility.

\paragraph{Oscillatory extrema count $N_{\text{extrema}}$.}
The number of local minima plus local maxima (using first-order neighbours)
in the non-interpolated HR residual sequence within each window.
This metric directly counts the number of reversal events,
complementing~$\chi(0)$ as an oscillatory complexity indicator.

\subsection{Sample entropy comparator metric}
\label{sec:sampen-method}

Sample entropy (SampEn) was computed as a comparator complexity metric.
For each 12-hour HR residual window, SampEn with template length $m=1$
was calculated using tolerance $r = 0.5 \times \sigma_{48}$, where
$\sigma_{48}$ is the standard deviation of the patient's full 48-hour
HR residual series, pre-computed once per patient and held constant
across all windows.
This patient-level reference follows the recommendation of Richman and
Moorman \cite{richman2000} to anchor $r$ to global signal variability.
Windows with fewer than~8 non-interpolated observations were excluded
(NaN rate 1.1\% at $r=0.5$; the standard $r=0.2$ yields 15.2\% NaN,
because the short 12-point window produces zero length-2 template
matches [$A=0$]).
An $m=2$ sensitivity analysis is reported in Supplementary Table~S1.

\subsection{External validation cohort (eICU-CRD)}
\label{sec:eicu-method}

The eICU Collaborative Research Database v2.0 (eICU-CRD)
\cite{pollard2018} was accessed under PhysioNet Data Use Agreement.
eICU-CRD covers 335 US ICUs participating in a telehealth critical care
programme.
The identical operational septic shock definition and $T_0$ assignment
were applied, yielding a 1:1 matched cohort (SOFA\,$\pm$2 risk-set
matching, seed 114514) of 587 matched pairs (587 shock, 587 control).
Time-series preprocessing and topological indicator computation
($\chi_{\text{HR}}(0)$, $N_{\text{extrema}}$, SampEn) followed the
same pipeline as MIMIC-IV.
The early window was defined as [$-$24,\,$-$12)\,h before $T_0$
(median $T_0$ observation window in eICU-CRD is approximately 24\,h,
precluding the [$-$48,\,$-$24)\,h definition used in MIMIC-IV);
the late window ($-12$ to $<$0\,h) was identical.
To account for the short and variable $T_0$ observation windows in
eICU-CRD, a T0-stratified sensitivity analysis was conducted:
associations were estimated separately in subsets with $\geq$24,
$\geq$48, and $\geq$72\,h of HR observations before $T_0$, since
signal accumulation requires sufficient observation time.
This analysis is post-hoc and exploratory.

\subsection{Statistical analysis}
\label{sec:stats}

For MIMIC-IV, patient-level summaries were computed as window means
within the \emph{early} window (window-centre $-48$ to $-24$\,h before
$T_0$) and \emph{late} window ($-12$ to $<0$\,h before $T_0$), and their
within-patient late-minus-early difference ($\Delta$).
For eICU-CRD, the early window was $-24$ to $-12$\,h (as specified
in Section~\ref{sec:eicu-method}); the late window was identical.

\paragraph{Between-group comparisons.}
Gaussian generalised estimating equations (GEE) \cite{liang1986} with
exchangeable within-cluster correlation and \texttt{stay\_id} clustering
were used for early, late, and $\Delta$ summaries.
All comparisons were subjected to global Holm--Bonferroni correction across
all metrics and time windows simultaneously ($6 \times 3 = 18$ tests total).

\paragraph{Temporal trend.}
Linear mixed-effects models (LMM) with form
$\texttt{indicator} \sim \texttt{time} \times \texttt{group}
+ (1 + \texttt{time} \mid \texttt{pair})$
were fitted over the full 48-hour window (time centred at $-24$\,h).
The time-by-group interaction coefficient~$\beta$ quantifies divergence
in trajectories between shock and control groups.

\paragraph{Conditional logistic regression.}
To estimate adjusted associations, conditional logistic regression stratified
by \texttt{matched\_pair\_id} was fitted with predictors:
(i)~late-window $\chi_{\text{HR}}(0)$ mean (primary predictor);
(ii)~invasive mechanical ventilation at the late-window start ($-12$\,h);
(iii)~sedative exposure in the prior 12\,h;
(iv)~beta-blocker exposure in the prior 12\,h;
(v)~ICU type (MICU reference); and
(vi)~monitoring modality (NBP reference).
Admission SOFA and time-to-event are structurally conditioned out by the
matched-pair likelihood.
Results are reported as odds ratios (OR) with 95\% confidence intervals.

\paragraph{Sensitivity analyses.}
\begin{enumerate}
  \item[\textbf{S1}]~Additional adjustment for late-window raw mean HR
    (key test: whether association survives mean-level confounding).
  \item[\textbf{S2}]~Additional adjustment for late-window raw mean HR
    and MAP simultaneously.
  \item[\textbf{S3}]~Early-window $\chi_{\text{HR}}(0)$ as predictor,
    adjusted for early-window raw mean HR and MAP
    (specification aligned with the AC1 early-window model in Fang et al.).
  \item[\textbf{S4}]~$N_{\text{extrema}}$ as primary predictor (base model
    and with additional late-window HR mean adjustment).
  \item[\textbf{S5}]~Restriction to pairs without sedative exposure in the
    prior 12\,h (sedation covariate omitted).
\end{enumerate}

\paragraph{Comparator analyses.}
Late-window HR~AC1 was analysed under the same conditional logistic
framework (M1, S1, S2, S3, S5), reproducing the analysis of Fang
et al.\ \cite{fang2026} within the same cohort.
SampEn ($m=1$, $r=0.5 \times \sigma_{48}$) was similarly analysed
under M1, S1, and S2.
The slightly smaller AC1 sample ($N=2{,}059$ in M1 versus
$N=2{,}156$ for Euler models) reflects NaN returns when within-window
residual variance is insufficient for a well-defined lag-1 correlation,
not a difference in the minimum-observation criterion.
Comparisons are therefore across a shared analytical framework but
not a strictly identical sample.

All analyses were performed in Python 3.14 using \texttt{pandas},
\texttt{numpy}, \texttt{scipy}, \texttt{statsmodels}, and \texttt{seaborn}.
Code is available at \url{https://github.com/} [to be added upon publication].

% ================================================================
\section{Results}
\label{sec:results}

\subsection{Cohort characteristics}
The MIMIC-IV discovery cohort comprised 2,172 analysable (stay,\,$T_0$)
observations (725 shock, 1,447 control) from 725 matched strata.
The shock group had higher in-hospital mortality (56.8\% vs.\ 21.3\%,
$p<0.001$) and a greater proportion with prior vasopressor exposure
before $T_0$ (93.2\% vs.\ 41.4\%, $p<0.001$), as expected.
Median age (65 vs.\ 63 years), sex (56.7\% vs.\ 59.6\% male), and
admission SOFA (6 [4--9] vs.\ 6 [4--8]) were comparable across matched
strata.
The eICU-CRD validation cohort comprised 587 matched pairs (587 shock,
587 control); full baseline characteristics for both cohorts are reported
in the companion paper \cite{fang2026}.

\subsection{Between-group comparisons of topological indicators}
\label{sec:res-gee}

Table~\ref{tab:gee} and Figure~\ref{fig:timeseries} summarise between-group
comparisons after global Holm correction.

\paragraph{HR Euler characteristic $\chi_{\text{HR}}(0)$.}
In the early window, $\chi_{\text{HR}}(0)$ was already significantly lower
in shock patients (mean 1.75 vs.\ 1.86; Holm $p=0.020$).
The difference was larger in the late window (1.62 vs.\ 1.89;
Holm $p<0.001$), and the within-patient late-minus-early change was also
significantly more negative in shock ($-0.131$ vs.\ $+0.030$;
Holm $p=0.013$).
The LMM time-by-group interaction was negative and significant
($\beta=-0.0053$ per hour, $p=0.006$), indicating a progressive divergence
between groups over the 48-hour window.

\paragraph{HR extrema count $N_{\text{extrema}}$.}
Late-window $N_{\text{extrema}}$ was significantly lower in shock
(5.16 vs.\ 5.53; Holm $p<0.001$), with a significant late-minus-early
change ($p<0.001$) and LMM interaction ($\beta=-0.0090$, $p=0.004$).
The early-window difference did not reach significance (Holm $p=1.00$).

\paragraph{MAP indicators.}
MAP~$\chi(0)$ showed no significant difference at any time point (all Holm
$p\geq0.574$).
Normalised total variation for both HR and MAP showed no consistent group
differences (all Holm $p>0.15$).

\subsection{Conditional logistic regression}
\label{sec:res-logit}

Table~\ref{tab:logit} presents the conditional logistic results.

\paragraph{$N_{\text{extrema}}$ (primary topological indicator).}
Lower late-window $N_{\text{extrema}}$ was associated with shock in
all models and was completely stable across mean-adjustment steps:
base model OR\,=\,0.88 (95\%\,CI 0.82--0.94), $p<0.001$
($N=2{,}156$; 720 informative strata);
additionally adjusted for late-window raw mean HR,
OR\,=\,0.88 (0.82--0.94), $p<0.001$;
additionally adjusted for both late-window raw mean HR and MAP,
OR\,=\,0.87 (0.81--0.94), $p<0.001$.
In the no-sedation subgroup (S5), the association strengthened:
OR\,=\,0.79 (0.67--0.93), $p=0.004$ ($N=544$; 218 strata).

\paragraph{$\chi_{\text{HR}}(0)$ (secondary topological indicator).}
Lower late-window $\chi_{\text{HR}}(0)$ was also associated with shock:
base model OR\,=\,0.75 (0.66--0.84), $p<0.001$;
after adjustment for late-window raw mean HR, OR\,=\,0.79 (0.70--0.89),
$p<0.001$;
after further adjustment for raw mean MAP, OR\,=\,0.86 (0.76--0.98),
$p=0.025$.
Early-window $\chi_{\text{HR}}(0)$ (S3, adjusted for early-window mean
HR and MAP) was similarly associated: OR\,=\,0.85 (0.74--0.98),
$p=0.027$.
In the no-sedation subgroup, OR\,=\,0.72 (0.58--0.89), $p=0.002$
($N=544$; 218 strata).

\subsection{Comparison with lag-1 autocorrelation and sample entropy}
\label{sec:res-comparators}

Figure~\ref{fig:sampenforest} shows the four-indicator comparison across
M1, S1, and S2 models.

\paragraph{Lag-1 autocorrelation (AC1).}
Under the identical conditional logistic framework, higher late-window
HR~AC1 was associated with shock in the base and single-covariate models
(M1: OR\,=\,1.60 [1.15--2.22], $p=0.005$; S1: OR\,=\,1.49 [1.06--2.08],
$p=0.021$), but lost significance after additional MAP adjustment
(S2: OR\,=\,1.23 [0.85--1.78], $p=0.271$) and in the no-sedation
subgroup (S5: OR\,=\,1.07 [0.59--1.95], $p=0.815$).
The key contrast is in S2: $N_{\text{extrema}}$ retained significance
(OR\,=\,0.87, $p<0.001$) and $\chi_{\text{HR}}(0)$ retained significance
(OR\,=\,0.86, $p=0.025$), while AC1 did not.

\paragraph{Sample entropy.}
SampEn ($m=1$, $r=0.5\times\sigma_{48}$) was not significantly associated
with shock in any model: M1 OR\,=\,0.86 (0.69--1.07), $p=0.178$;
S1 OR\,=\,0.81 (0.64--1.02), $p=0.067$;
S2 OR\,=\,0.87 (0.68--1.12), $p=0.284$.
Non-significance was not solely attributable to reduced sample size:
SampEn M1 ($N=2{,}122$) was comparable to the Euler M1 sample
($N=2{,}156$), but SampEn confidence intervals were wider, reflecting
greater between-subject variability.
At the standard tolerance ($r=0.2\times\sigma_{48}$), 15.2\% of
late-window SampEn values were undefined (versus 0.3\% for
$N_{\text{extrema}}$), because the 12-point window yields zero
length-2 template matches ($A=0$); relaxing to $r=0.5$ reduced the
NaN rate to 1.1\% but did not produce significant associations.

\subsection{External validation in eICU-CRD}
\label{sec:res-eicu}

Table~\ref{tab:eicu} shows the T0-stratified external validation results
in eICU-CRD (587 matched pairs).
In the full cohort, no indicator reached significance, consistent with
the short median $T_0$ observation window ($\approx$24\,h) in eICU-CRD
that limits early-window signal accumulation.
In the subset with $T_0 \geq 72$\,h before vasopressor initiation
(103 pairs, 100 informative strata), $N_{\text{extrema}}$ was
significantly associated with shock (OR\,=\,0.80, 95\%\,CI 0.65--0.98;
$p=0.034$), closely replicating the MIMIC-IV estimate (OR\,=\,0.88;
$p<0.001$).
$\chi_{\text{HR}}(0)$ showed a consistent direction (OR\,=\,0.72,
0.50--1.04; $p=0.084$) but did not reach the pre-specified $\alpha=0.05$
threshold.
SampEn showed no consistent direction and was non-significant at all
$T_0$ thresholds (OR\,=\,1.17 at $T_0\geq 72$\,h; $p=0.625$).
AC1, although numerically elevated at $T_0\geq 72$\,h
(OR\,=\,2.68, 0.88--8.12; $p=0.082$), had wide confidence intervals
reflecting the small stratum count.

% ================================================================
\section{Discussion}
\label{sec:discussion}

In two independent ICU cohorts, reduced topological complexity of hourly
heart rate---specifically the local-extrema count $N_{\text{extrema}}$---was
associated with late-onset septic shock before vasopressor initiation.
$N_{\text{extrema}}$ was the most robust indicator: its association in
MIMIC-IV (OR\,=\,0.88; $p<0.001$) was completely unchanged after adjusting
for concurrent mean HR and MAP, and it replicated in eICU-CRD at
$T_0\geq 72$\,h (OR\,=\,0.80; $p=0.034$).
The Euler characteristic $\chi_{\text{HR}}(0)$ showed a stronger unadjusted
association (OR\,=\,0.75) that partially attenuated with full mean adjustment
(OR\,=\,0.86), indicating some overlap with mean-level deterioration.
To our knowledge, this is the first application of sublevel-set Euler
characteristic analysis to ICU vital-sign early-warning signals; prior TDA
work in medicine has focused on persistent homology of ECG waveforms
\cite{chung2021} rather than hourly EHR time series.

The physiological substrate is progressive failure of autonomic cardiovascular
regulation in evolving sepsis \cite{decastilho2018,carrara2021}.
Under normal conditions, HR oscillates bidirectionally around its local trend
through sympathetic--parasympathetic interplay; $N_{\text{extrema}}$ and
$\chi_{\text{HR}}(0)$ directly index this bidirectionality.
As septic physiology evolves---through autonomic neuropathy, catecholamine
desensitisation, and microcirculatory failure---HR loses these reversals and
becomes persistently elevated with fewer oscillatory crossings.
Critically, this structural loss is separable from simple tachycardia:
$N_{\text{extrema}}$ was unchanged after adjusting for mean HR (OR\,=\,0.88,
$p<0.001$), and it remained significant in patients not receiving sedatives
(OR\,=\,0.79, $p=0.004$), arguing against sedation confounding as the
primary driver.
This finding is consistent with the broader HRV literature showing that
complexity measures provide prognostic information beyond the HR level
in critical illness \cite{moorman2011}.

The differential behaviour of AC1 and topological indicators under mean
adjustment illuminates a key limitation of the critical slowing down
framework as applied to clinical vital signs.
AC1 measures autocorrelation persistence---a property that co-varies with
the process mean when sustained tachycardia is the dominant signal: HR
remaining above its detrended mean for extended runs inflates AC1.
After adjusting for this mean shift (S2 model), AC1 lost significance
(OR\,=\,1.23; $p=0.271$), and was non-significant in the no-sedation
subgroup and in eICU-CRD, suggesting its apparent signal was largely
a proxy for overt haemodynamic deterioration rather than an independent
structural feature.
$N_{\text{extrema}}$ and $\chi_{\text{HR}}(0)$, by contrast, are computed
relative to the current trend and therefore remained significant after
empirical adjustment for the mean level in our data.

Sample entropy also failed to capture an independent association with
shock (OR\,=\,0.86; $p=0.178$ in the unadjusted model), and showed no
consistent direction in eICU-CRD (OR\,=\,1.17 at $T_0\geq 72$\,h).
This failure reflects a structural limitation of standard entropy metrics
when applied to sparse EHR time series: at the conventional tolerance
($r=0.2\times$local SD), 15.2\% of 12-hour windows produced undefined
SampEn values because no length-2 template matches existed within the
approximately 12 available observations.
Relaxing $r$ to $0.5\times\sigma_{48}$ reduced NaN rates to 1.1\% but
did not recover a significant association.
This highlights a practical advantage of topological indicators for routine
EHR monitoring: $N_{\text{extrema}}$ and $\chi_{\text{HR}}(0)$ require no
template parameters, are undefined in fewer than 0.5\% of windows, and are
interpretable from as few as six observations---conditions routinely met
in hourly EHR charting.

Several limitations should be noted.
This analysis is exploratory and hypothesis-generating; neither topological
indicator had a pre-registered primary hypothesis, and the analysis was
conducted post-hoc within an existing cohort.
MIMIC-IV is a single academic centre, and eICU-CRD, while multi-centre,
represents primarily US academic hospitals with telehealth infrastructure;
generalisability to community hospitals and non-US settings requires further
study.
The early-window definition differs between cohorts (MIMIC-IV:
$-48$ to $-24$\,h; eICU: $-24$ to $-12$\,h), precluding direct comparison
of early-window estimates.
The hourly EHR sampling rate limits observations to approximately 6--12
per 12-hour window; waveform-resolution data could substantially improve
precision and enable multiscale analysis.
The T0-stratified eICU analysis is not a pre-specified primary endpoint
and should be interpreted with appropriate caution.
Notwithstanding, the consistent replication of $N_{\text{extrema}}$ in
eICU-CRD at $T_0\geq 72$\,h, its complete robustness to mean adjustment
across all sensitivity models, and the structural failure of SampEn on
sparse EHR data together provide a coherent argument for topological
indicators as the preferred framework for this application.
These findings support prospective evaluation of $N_{\text{extrema}}$ and
$\chi_{\text{HR}}(0)$ as complements to current ICU early-warning systems;
implementation requires no additional monitoring hardware, and the
$O(n)$ computation is compatible with real-time clinical deployment.

% ================================================================
\section*{Declarations}

\paragraph{Data availability.}
MIMIC-IV v3.1 and eICU-CRD v2.0 are available via PhysioNet
(\url{https://physionet.org}) under Data Use Agreements.
Derived (de-identified) intermediate files are available at
[GitHub link upon publication].

\paragraph{Code availability.}
All analysis code (Python scripts) is available at
[GitHub link upon publication].

\paragraph{Funding.}
No specific funding was received for this exploratory analysis.

\paragraph{Conflicts of interest.}
None declared.

\paragraph{Author contributions.}
C.F.: conceptualisation, data curation, formal analysis, methodology,
writing—original draft.
H.Y., Q.L.: supervision, writing—review \& editing.

% ================================================================
\section*{Abbreviations}
\begin{tabular}{ll}
  AC1         & Lag-1 autocorrelation \\
  CSD         & Critical slowing down \\
  EHR         & Electronic health record \\
  eICU-CRD    & eICU Collaborative Research Database \\
  EWS         & Early warning signal \\
  GEE         & Generalised estimating equations \\
  HR          & Heart rate \\
  HRV         & Heart rate variability \\
  ICU         & Intensive care unit \\
  LMM         & Linear mixed-effects model \\
  MAP         & Mean arterial pressure \\
  MSE         & Multiscale entropy \\
  OR          & Odds ratio \\
  SampEn      & Sample entropy \\
  SOFA        & Sequential Organ Failure Assessment \\
  TDA         & Topological data analysis \\
  $\chi(0)$   & Euler characteristic at threshold zero \\
\end{tabular}

% ================================================================
\clearpage
\appendix
\section*{Appendix: Mathematical Note on Euler Characteristic}
\label{app:math}

For a discrete one-dimensional signal
$r = (r_1, \ldots, r_n) \in \mathbb{R}^n$,
the \emph{sublevel set} at threshold $\lambda$ is
$S_\lambda = \{i \in \{1,\ldots,n\} : r_i \leq \lambda\}$,
viewed as a subset of the path graph on $n$ vertices
(i.e.\ two adjacent vertices $i$ and $i+1$ both in $S_\lambda$
are connected by an edge).
The Euler characteristic of $S_\lambda$ is:
\[
  \chi(S_\lambda)
  = \beta_0(S_\lambda) - \beta_1(S_\lambda)
  = \beta_0(S_\lambda),
\]
since a subgraph of a path graph contains no cycles ($\beta_1 = 0$).
Here $\beta_0(S_\lambda)$ denotes the number of connected components.
For a path graph, connected components of $S_\lambda$ correspond to
maximal contiguous runs of indices with $r_i \leq \lambda$.
Setting $\lambda = 0$ and evaluating equation~(\ref{eq:euler}) follows directly.

The \emph{Euler characteristic curve} $\lambda \mapsto \chi(S_\lambda)$
is a complete topological summary of $r$ under the sublevel-set filtration
\cite{edelsbrunner2010}.
At $\lambda = -\infty$, $\chi = 0$; at each local minimum $r_{\min}$,
$\chi$ increases by~1 (a new component is born); at each local maximum
$r_{\max}$, $\chi$ decreases by~1 (two components merge); at
$\lambda = +\infty$, $\chi = 1$ (one connected component).
Our scalar $\chi_{\text{HR}}(0)$ is simply the evaluation of this curve
at the physiologically meaningful threshold of zero (the detrended mean),
rendering it directly interpretable as the count of below-mean oscillatory
excursions.

% ================================================================
\clearpage

\begin{table}[htbp]
\centering
\caption{Between-group comparisons of topological indicators (GEE,
Gaussian family, exchangeable correlation, \texttt{stay\_id} clustering).
$p$-values are Holm--Bonferroni corrected globally across all 18 tests
(6 metrics $\times$ 3 time summaries).
Values are group means.}
\label{tab:gee}
\small
\resizebox{\textwidth}{!}{%
\begin{tabular}{lcccccccccc}
\toprule
& \multicolumn{3}{c}{\textbf{Early window}} & \multicolumn{3}{c}{\textbf{Late window}} & \multicolumn{3}{c}{\textbf{$\Delta$ (late$-$early)}} \\
\cmidrule(lr){2-4}\cmidrule(lr){5-7}\cmidrule(lr){8-10}
\textbf{Indicator} & Shock & Control & Holm $p$ & Shock & Control & Holm $p$ & Shock & Control & Holm $p$ \\
\midrule
HR $\chi(0)$        & 1.75 & 1.86 & 0.020   & 1.62 & 1.89 & $<$0.001 & $-$0.131 & $+$0.030 & 0.013 \\
HR $N_{\text{ext}}$ & 5.48 & 5.52 & 1.000   & 5.16 & 5.53 & $<$0.001 & $-$0.315 & $+$0.020 & $<$0.001 \\
MAP $\chi(0)$       & 2.39 & 2.35 & 1.000   & 2.28 & 2.33 & 1.000    & $-$0.107 & $-$0.017 & 0.574 \\
MAP $N_{\text{ext}}$& 6.06 & 5.98 & 0.964   & 5.73 & 5.92 & 0.055    & $-$0.338 & $-$0.066 & 0.012 \\
HR total var$^\dag$ & 55.1 & 57.9 & 0.562   & 56.1 & 57.8 & 1.000    & $+$1.01  & $-$0.04  & 1.000 \\
MAP total var$^\dag$& 78.4 & 82.9 & 0.150   & 80.9 & 83.4 & 1.000    & $+$2.52  & $+$0.52  & 1.000 \\
\midrule
LMM interaction & \multicolumn{3}{c}{HR $\chi(0)$: $\beta=-0.0053$, $p=0.006$} &
\multicolumn{3}{c}{HR $N_{\text{ext}}$: $\beta=-0.0090$, $p=0.004$} &
\multicolumn{3}{c}{MAP: all $p>0.05$} \\
\bottomrule
\end{tabular}}
\vspace{2pt}\\
\raggedright
\footnotesize $^\dag$ Normalised total variation: $\sum|\Delta r| / (n_{\text{actual}}-1)$.
\end{table}

\begin{table}[htbp]
\centering
\caption{Conditional logistic regression results (stratified by
\texttt{matched\_pair\_id}).
Primary predictor: late-window HR $\chi(0)$ mean (models M1--S3, S5)
or $N_{\text{extrema}}$ (models S4, S4b).
Comparator: late-window HR AC1 under identical specifications.
OR\,$<$1 for $\chi(0)$ and $N_{\text{extrema}}$ corresponds to lower
topological complexity associated with higher shock odds.}
\label{tab:logit}
\small
\begin{tabular}{llccc}
\toprule
\textbf{Model} & \textbf{Predictor} & \textbf{N} & \textbf{OR (95\% CI)} & \textbf{$p$} \\
\midrule
\multicolumn{5}{l}{\textit{HR Euler $\chi(0)$ — primary and sensitivity analyses}} \\
M1 Base model            & $\chi_{\text{HR}}(0)$ late       & 2,156 & 0.75 (0.66--0.84) & $<$0.001 \\
S1 + late raw HR mean    & $\chi_{\text{HR}}(0)$ late       & 2,156 & 0.79 (0.70--0.89) & $<$0.001 \\
S2 + late raw HR + MAP   & $\chi_{\text{HR}}(0)$ late       & 2,138 & 0.86 (0.76--0.98) & 0.025    \\
S3 Early window$^a$      & $\chi_{\text{HR}}(0)$ early      & 2,171 & 0.85 (0.74--0.98) & 0.027    \\
S4 $N_{\text{extrema}}$ base & $N_{\text{extrema}}$ late    & 2,156 & 0.88 (0.82--0.94) & $<$0.001 \\
S4b + late raw HR mean   & $N_{\text{extrema}}$ late        & 2,156 & 0.88 (0.82--0.94) & $<$0.001 \\
S5 No-sedation subgroup  & $\chi_{\text{HR}}(0)$ late       &   544 & 0.72 (0.58--0.89) & 0.002    \\
\midrule
\multicolumn{5}{l}{\textit{HR AC1 — comparator under identical framework}} \\
M1 Base model            & AC1 late           & 2,059 & 1.60 (1.15--2.22) & 0.005    \\
S1 + late raw HR mean    & AC1 late           & 2,059 & 1.49 (1.06--2.08) & 0.021    \\
S2 + late raw HR + MAP   & AC1 late           & 2,041 & 1.23 (0.85--1.78) & 0.271    \\
S3 Early window$^a$      & AC1 early          & 2,155 & 1.11 (0.73--1.69) & 0.618    \\
S5 No-sedation subgroup  & AC1 late           &   526 & 1.07 (0.59--1.95) & 0.815    \\
\bottomrule
\end{tabular}
\vspace{2pt}\\
\raggedright
\footnotesize $^a$ S3 additionally adjusts for early-window raw mean HR and MAP
(aligned with the early-window AC1 specification in Fang et al.\ \cite{fang2026}).
All models include: invasive ventilation at $-12$\,h, sedative exposure
(omitted in S5), beta-blocker exposure, ICU type, and monitoring modality.
\end{table}

\begin{sidewaystable}[htbp]
\centering
\caption{External validation: T0-stratified conditional logistic regression
in eICU-CRD (1:1 matched cohort).
$T_0$-stratified subsets restrict to pairs with $\geq H$\,h of HR
observations before vasopressor initiation.
MIMIC-IV row shown for reference (M1 model specification).
$T_0$-stratified analysis is post-hoc and exploratory.}
\label{tab:eicu}
\small
\begin{tabular}{lcccccccc}
\toprule
& \multicolumn{2}{c}{\textbf{$N_{\text{extrema}}$}}
& \multicolumn{2}{c}{\textbf{Euler $\chi(0)$}}
& \multicolumn{2}{c}{\textbf{SampEn}}
& \multicolumn{2}{c}{\textbf{AC1}} \\
\cmidrule(lr){2-3}\cmidrule(lr){4-5}\cmidrule(lr){6-7}\cmidrule(lr){8-9}
\textbf{Subset ($N$ pairs)} & OR (95\%CI) & $p$ & OR (95\%CI) & $p$
& OR (95\%CI) & $p$ & OR (95\%CI) & $p$ \\
\midrule
\multicolumn{9}{l}{\textit{eICU-CRD}} \\
Full cohort (587)  & 0.94 (0.86--1.02) & 0.142 & 0.97 (0.81--1.16) & 0.756 & 1.00 (0.76--1.31) & 0.978 & 1.21 (0.75--1.97) & 0.435 \\
$T_0\geq 24$\,h (297) & 0.96 (0.86--1.07) & 0.455 & 1.02 (0.81--1.29) & 0.852 & 1.04 (0.72--1.52) & 0.819 & 1.50 (0.78--2.90) & 0.224 \\
$T_0\geq 48$\,h (154) & 0.92 (0.78--1.07) & 0.281 & 0.84 (0.62--1.14) & 0.262 & 1.14 (0.67--1.96) & 0.628 & 1.21 (0.50--2.93) & 0.673 \\
$T_0\geq 72$\,h (103)$^\star$ & \textbf{0.80 (0.65--0.98)} & \textbf{0.034} & 0.72 (0.50--1.04) & 0.084 & 1.17 (0.62--2.21) & 0.625 & 2.68 (0.88--8.12) & 0.082 \\
\midrule
\multicolumn{9}{l}{\textit{MIMIC-IV (reference)}} \\
Full cohort (725) & 0.88 (0.82--0.94) & $<$0.001 & 0.75 (0.66--0.84) & $<$0.001 & 0.86 (0.69--1.07) & 0.178 & 1.60 (1.15--2.22) & 0.005 \\
\bottomrule
\end{tabular}
\vspace{2pt}\\
\raggedright
\footnotesize $^\star$ T0-stratified analysis is post-hoc exploratory; interpret with caution.
eICU-CRD covariates: invasive ventilation, sedative exposure, beta-blocker exposure,
and arterial line monitoring (ICU type absorbed by 1:1 matching and excluded).
SampEn: $m=1$, $r=0.5\times\sigma_{48}$. AC1 eICU sample slightly smaller
(full cohort $N=1{,}102$) due to NaN from insufficient within-window variance.
\end{sidewaystable}

% ================================================================
\clearpage
\begin{figure}[htbp]
  \centering
  \includegraphics[width=\textwidth]{%
    ../output/fig_euler_timeseries.png}
  \caption{Time series of six topological indicators from $-48$ to $0$\,h
  relative to $T_0$ (3 rows $\times$ 2 columns).
  Mean~$\pm$~95\% confidence interval (shaded); LOESS-smoothed trend
  (solid line).
  Green dotted line: early/late window boundary ($-24$\,h);
  orange dotted line: late window start ($-12$\,h);
  grey dashed line: $T_0$.
  \textbf{Row 1 (left/right):} MAP Euler $\chi(0)$ and HR Euler $\chi_{\text{HR}}(0)$.
  \textbf{Row 2:} MAP $N_{\text{extrema}}$ and HR $N_{\text{extrema}}$.
  \textbf{Row 3:} MAP total variation and HR total variation.
  HR Euler $\chi_{\text{HR}}(0)$ and HR $N_{\text{extrema}}$ decline
  progressively in the shock group relative to controls from approximately
  $-24$\,h onward; MAP indicators and total variation show no consistent
  group divergence.}
  \label{fig:timeseries}
\end{figure}

\begin{figure}[htbp]
  \centering
  \includegraphics[width=\textwidth]{%
    ../output/fig_euler_vs_ac1_forest.png}
  \caption{Forest plot comparing OR (95\%\,CI) for HR Euler $\chi(0)$
  (left panel) and HR AC1 (right panel) across corresponding model
  specifications.
  Red symbols indicate $p<0.05$; grey symbols indicate $p\geq0.05$.
  The key difference is in S2: $\chi(0)$ retains significance after
  adjustment for both late-window raw mean HR and MAP,
  whereas AC1 does not.}
  \label{fig:forest}
\end{figure}

\begin{figure}[htbp]
  \centering
  \includegraphics[width=\textwidth]{%
    ../output/fig_sampen_4way_forest.png}
  \caption{Four-indicator comparison across M1, S1, and S2 model
  specifications (MIMIC-IV discovery cohort).
  Panels show OR (95\%\,CI) for HR $N_{\text{extrema}}$ (top left),
  HR Euler $\chi(0)$ (top right), SampEn $m=1$, $r=0.5\sigma_{48}$
  (bottom left), and HR AC1 (bottom right).
  Red symbols indicate $p<0.05$; grey indicate $p\geq0.05$.
  $N_{\text{extrema}}$ is the only indicator that remains significant
  across all three model specifications.}
  \label{fig:sampenforest}
\end{figure}

\begin{figure}[htbp]
  \centering
  \includegraphics[width=0.9\textwidth]{%
    ../output/fig_euler_delta.png}
  \caption{Early-to-late change ($\Delta$) in topological indicators
  by group.
  Each panel shows late-minus-early difference for one indicator;
  horizontal dashed line at~0.
  Both HR $\chi(0)$ and HR $N_{\text{extrema}}$ show a significantly more
  negative $\Delta$ in the shock group (Holm $p\leq0.013$), indicating
  progressive loss of oscillatory complexity approaching shock onset.}
  \label{fig:delta}
\end{figure}

% ================================================================
\bibliographystyle{unsrtnat}
\begin{thebibliography}{99}

\bibitem{singer2016}
Singer M, Deutschman CS, Seymour CW, et al.
The Third International Consensus Definitions for Sepsis and Septic Shock
(Sepsis-3).
\textit{JAMA}. 2016;315(8):801--810.

\bibitem{scheffer2009}
Scheffer M, Bascompte J, Brock WA, et al.
Early-warning signals for critical transitions.
\textit{Nature}. 2009;461(7260):53--59.

\bibitem{dakos2012}
Dakos V, Carpenter SR, Brock WA, et al.
Methods for detecting early warnings of critical transitions in time series
illustrated using simulated ecological data.
\textit{PLOS ONE}. 2012;7(7):e41010.

\bibitem{olderikkert2016}
Olde Rikkert MG, Dakos V, Buchman TG, et al.
Slowing down of recovery as generic risk marker for acute severity transitions
in chronic diseases.
\textit{Crit Care Med}. 2016;44(3):601--606.

\bibitem{khalil2025}
Khalil L, Cliff ERS, Marhoon N, et al.
Transitions in intensive care: Investigating critical slowing down
post extubation.
\textit{PLOS ONE}. 2025;20(1):e0316908.
DOI: 10.1371/journal.pone.0316908

\bibitem{decastilho2018}
de Castilho FM, Ribeiro ALP, da Silva JLP, Nobre V, de Sousa MR.
Heart rate variability as predictor of mortality in sepsis: a prospective
cohort study.
\textit{PLOS ONE}. 2018;13(9):e0203487.

\bibitem{carrara2021}
Carrara M, Laffey JG, Pham T, et al.
Heart rate variability analysis in critically ill patients: a systematic
review of methodological and clinical issues.
\textit{J Clin Monit Comput}. 2021;35(3):619--640.
DOI: 10.1007/s10877-020-00497-3

\bibitem{moorman2011}
Moorman JR, Carlo WA, Kattwinkel J, et al.
Mortality reduction by heart rate characteristic monitoring in very low
birthweight neonates: a randomized trial.
\textit{J Pediatr}. 2011;159(6):900--906.

\bibitem{lake2002}
Lake DE, Richman JS, Griffin MP, Moorman JR.
Sample entropy analysis of neonatal heart rate variability.
\textit{Am J Physiol Regul Integr Comp Physiol}. 2002;283(3):R789--R797.

\bibitem{bravi2012}
Bravi A, Longtin A, Seely AJE.
Review and classification of variability analysis techniques with clinical
applications.
\textit{Biomed Eng Online}. 2012;11:90.

\bibitem{edelsbrunner2010}
Edelsbrunner H, Harer J.
\textit{Computational Topology: An Introduction}.
American Mathematical Society; 2010.

\bibitem{chung2021}
Chung YM, Hu CS, Lo YL, Wu HT.
A persistent homology approach to heart rate variability analysis with
an application to sleep-wake classification.
\textit{Front Physiol}. 2021;12:637684.

\bibitem{johnson2023}
Johnson AEW, Bulgarelli L, Shen L, et al.
MIMIC-IV, a freely accessible electronic health record dataset.
\textit{Sci Data}. 2023;10:1.

\bibitem{fang2026}
Fang C, Yu H, Luo Q.
Heart rate lag-1 autocorrelation is associated with pre-shock dynamics
in late-onset septic shock: a single-cohort exploratory analysis of MIMIC-IV.
\textit{[Journal]}. 2026 [under review].

\bibitem{liang1986}
Liang KY, Zeger SL.
Longitudinal data analysis using generalized linear models.
\textit{Biometrika}. 1986;73(1):13--22.

\bibitem{pollard2018}
Pollard TJ, Johnson AEW, Raffa JD, Celi LA, Mark RG, Badawi O.
The eICU Collaborative Research Database, a freely available
multi-center database for critical care research.
\textit{Sci Data}. 2018;5:180178.
DOI: 10.1038/sdata.2018.178

\bibitem{richman2000}
Richman JS, Moorman JR.
Physiological time-series analysis using approximate entropy and sample
entropy.
\textit{Am J Physiol Heart Circ Physiol}. 2000;278(6):H2039--H2049.

\end{thebibliography}

\end{document}
